\documentclass[11pt,a4paper]{article}

\usepackage[margin=2.3cm]{geometry}
\usepackage{amsmath,amssymb,amsthm}
\usepackage[T1]{fontenc}
\usepackage{tgschola}
\usepackage{fouriernc}
\usepackage{microtype}
\usepackage{tfrupee}
\linespread{1.04}
\usepackage{xcolor}
\usepackage{titlesec}
\usepackage{tikz}
\usetikzlibrary{angles,quotes,calc,arrows.meta}
\usepackage{enumitem}
\usepackage{array}
\usepackage{needspace}
\usepackage{fancyhdr}
\usepackage[hidelinks]{hyperref}

\definecolor{accent}{RGB}{38,52,70}
\definecolor{panel}{RGB}{236,238,241}

\newtheoremstyle{brand}{9pt}{9pt}{\itshape}{}{\bfseries\color{accent}}{.}{ }{}
\newtheoremstyle{branddef}{9pt}{9pt}{}{}{\bfseries\color{accent}}{.}{ }{}
\theoremstyle{brand}
\newtheorem{theorem}{Theorem}
\newtheorem{corollary}{Corollary}
\theoremstyle{branddef}
\newtheorem{definition}{Definition}
\newtheorem{example}{Example}[section]
\newtheorem{illus}{Illustration}
\newtheorem{remark}{Remark}
\newtheorem*{activity}{Activity}

\newcommand{\tr}{\triangle}
\newenvironment{exercise}[1]
  {\par\medskip\Needspace{7\baselineskip}\noindent{\large\bfseries\color{accent}Exercise #1}
   \par\smallskip
   \begin{enumerate}[leftmargin=*,itemsep=4pt]}
  {\end{enumerate}}
\definecolor{muted}{RGB}{70,72,76}
\newcommand{\nl}{\\[3pt]}
\ifteacher
\newenvironment{answerblock}
  {\par\smallskip\Needspace{10\baselineskip}\begingroup\small\linespread{1.15}\selectfont\color{muted}\noindent\textit{Answers.}\par}
  {\par\endgroup\medskip}
\newcommand{\ans}[1]{\par\smallskip\Needspace{9\baselineskip}\noindent\textbf{#1}\ }
\else
\usepackage{comment}
\excludecomment{answerblock}
\fi
\newcommand{\figph}[1]{\begin{center}\fbox{\parbox{.85\linewidth}{\small FIGURE: #1}}\end{center}}
% right-angle mark at vertex #1 between the directions to #2 and #3
\newcommand{\ramark}[3]{\coordinate (ra1) at ($(#1)!2mm!(#2)$);\coordinate (ra2) at ($(#1)!2mm!(#3)$);\draw[thin] (ra1) -- ($(ra1)+(ra2)-(#1)$) -- (ra2);}
% side label: #4 placed beside segment #2--#3, on the right-hand side when walking from #2 to #3
% (for a triangle named anticlockwise this is the outside); optional #1 = position along the segment
\newcommand{\slab}[4][0.5]{\path (#2) -- (#3) node[pos=#1, auto=right, inner sep=2.5pt] {#4};}
% dimension line for a whole side: drawn #3 away from segment #1--#2 on its right-hand side, labelled #4
\newcommand{\dimlab}[4]{\draw[thin,{Bar[width=5pt]}-{Bar[width=5pt]}] ($(#1)!#3!-90:(#2)$) -- ($(#2)!#3!90:(#1)$) node[midway, auto=right, inner sep=2pt] {#4};}
\newcommand{\exfig}[1]{\begin{center}\begin{tikzpicture}[scale=.75, every node/.style={font=\small}]#1\end{tikzpicture}\end{center}}
% rough answer diagram (teacher copy only; it sits inside answerblock)
\newcommand{\ansfig}[1]{\begin{center}\begin{tikzpicture}[scale=.6, every node/.style={font=\footnotesize}, draw=muted]#1\end{tikzpicture}\end{center}}
\newcommand{\rememberbox}[1]{%
  \par\medskip\noindent
  \colorbox{panel}{%
    \parbox{\dimexpr\linewidth-2\fboxsep\relax}{%
      \medskip
      \textbf{\color{accent}\large Things to Remember}
      \par\smallskip #1 \medskip}}
  \par\medskip}
\ifteacher
\newcommand{\facilitatornote}[1]{%
  \par\medskip\noindent
  \colorbox{panel}{%
    \parbox{\dimexpr\linewidth-2\fboxsep\relax}{%
      \smallskip
      {\small\textbf{Note for the teacher.}\ #1}
      \smallskip}}
  \par\medskip}
\else
\newcommand{\facilitatornote}[1]{}
\fi

\newcommand{\ModuleTitle}{CIRCLES}
\newcommand{\PartTitle}{Lesson 1}
\newcommand{\ModuleDescriptor}{Grade 9 mathematics facilitator copy}

\pagestyle{plain}
\setlength{\parskip}{4pt}
\setlength{\parindent}{0pt}
% keep tall fractions in running text from touching the line above or below
\setlength{\lineskiplimit}{2pt}
\setlength{\lineskip}{3pt}


\titleformat{\section}{\large\bfseries}{\thesection}{0.8em}{}
\titlespacing*{\section}{0pt}{16pt}{6pt}
\titleformat{\subsection}{\normalsize\bfseries}{\thesubsection}{0.6em}{}
% plain ruled panels instead of coloured boxes
\renewcommand{\colorbox}[2]{\par\smallskip\noindent\textcolor{black!35}{\rule{\linewidth}{0.4pt}}\par\nointerlineskip\vspace{3pt}\noindent#2\par\vspace{2pt}\noindent\textcolor{black!35}{\rule{\linewidth}{0.4pt}}\par\smallskip}
